% ================================================================
\section{The Smith--Ramanujan Decomposition}
\label{sec:smith}
% ================================================================

The Smith--Ramanujan subsystem (\lean{Physics/GramWiring/*.lean},
\lean{Physics/Mertens/*.lean}, 18 files, zero sorry) decomposes
the Gram matrix through the Smith Normal Form, revealing the
Jordan totient eigenvalues and establishing an unconditional
spectral gap for the Ramanujan quadratic form.

\subsection{The Ramanujan Matrix}

The Ramanujan matrix $R$ with entries
\begin{equation}
  R(j,k) = \frac{\gcd(j,k)^2}{jk} - \frac{1}{jk}
  \label{eq:ramanujan-entry}
\end{equation}
is the ``arithmetic part'' of the Gram matrix. The Cathedral proves
(in \lean{RamanujanBridge.lean}) that the full Gram entry decomposes as:
\begin{equation}
  G(j,k) = R(j,k) + \text{(analytic correction)}
  \label{eq:gram-ramanujan}
\end{equation}
The Ramanujan matrix captures the \emph{algebraic} interaction between
modes $j$ and $k$ through their greatest common divisor, while the
analytic correction involves transcendental quantities
($\ln 2\pi$, $\gamma$, cotangent sums).

\subsection{The Smith Normal Form and Sum of Squares}

The Smith witness (\lean{SmithWitness.lean}, 50K bytes) proves that
the Ramanujan quadratic form has a sum-of-squares decomposition:
\begin{equation}
  v^T R\, v = \frac{1}{12} \sum_{d=1}^{N} J_2(d)\, y_d^2
  \label{eq:sos}
\end{equation}
where $J_2(d) = d^2 \prod_{p | d}(1 - 1/p^2)$ is the
\emph{Jordan totient function} and $y_d$ are linear combinations of the
$v_k$'s defined by the Smith decomposition.

\begin{correspondence}[Smith Normal Form = Spectral Diagonalization]
The Smith Normal Form of the $\gcd^2$ matrix factorizes as
$R = D^{-1} \Phi\, J\, \Phi^T D^{-1}/12$, where $\Phi$ is the
Euler totient matrix, $J$ is the diagonal matrix of Jordan
totients, and $D$ is the diagonal normalization. This is
the \emph{spectral diagonalization} of the Ramanujan scattering
matrix: the Jordan totients $J_2(d)$ are the eigenvalues,
and the Euler totient matrix~$\Phi$ provides the change of basis.

The physical content:
\begin{itemize}
  \item $J_2(d) > 0$ for all $d$ (since $1 - 1/p^2 > 0$), so
    $v^T R\, v \geq 0$---the Ramanujan matrix is
    \textbf{positive semidefinite}. The scattering matrix is
    ``reflection positive.''
  \item Each term $J_2(d) y_d^2$ represents the energy in the
    $d$-th ``stratum''---modes whose arithmetic structure is
    controlled by the divisor~$d$.
  \item The SOS structure implies that cancellation in $v^T R\, v$
    requires $y_d = 0$ for all~$d$---a strong constraint on the
    weight vector.
\end{itemize}
\end{correspondence}

\subsection{The East--West Principle}

The Smith--Franel Bridge (\lean{SmithFranelBridge.lean}) establishes
a profound asymmetry between two completeness problems:

\begin{principle}[The East--West Principle]
\begin{itemize}
  \item \textbf{The East Wing (unconditional):} The functions
    $\{kt\}$ for $k = 1, 2, \ldots$ are complete in $L^2(0,1)$.
    This is the Franel--Landau completeness theorem, related to
    Farey fractions. The Ramanujan quadratic form
    $\sigma(N) = \inf_v v^T R\, v \to 0$ unconditionally
    (proved from Euclid's theorem in \lean{SmithSpectralGap.lean}).

  \item \textbf{The West Wing (conditional):} The functions
    $\{1/(kx)\}$ for $k = 2, 3, \ldots$ are complete in $L^2(0,1)$
    \emph{if and only if} the Riemann Hypothesis holds. This is the
    Nyman--Beurling theorem.
\end{itemize}

The East Wing uses the ``Ramanujan'' inner product
$\langle \{jt\}, \{kt\} \rangle = R(j,k)$, which has no transcendental
content. The West Wing uses the ``Gram'' inner product
$\langle \{1/(jx)\}, \{1/(kx)\} \rangle = G(j,k)$, which
involves the Mellin-sensitive cotangent sums.

Both wings share the same underlying arithmetic (the GCD structure),
but they differ in their \emph{analytic weight}: the East Wing
sees only algebraic correlations (GCD power), while the West Wing
sees the full transcendental correlator (cotangent phase shifts).
The Riemann Hypothesis is the assertion that the transcendental
corrections do not obstruct completeness---that the algebra
determines the analysis.
\end{principle}

\subsection{The Spectral Gap: $\sigma(N) \to \infty$}

The Smith spectral gap (\lean{SmithSpectralGap.lean}) proves the key
unconditional result:
\begin{equation}
  \sigma(N) = \sum_{k=1}^{N} J_2(k) \to \infty
  \quad\text{as } N \to \infty
  \label{eq:sigma-diverges}
\end{equation}
This follows from Euclid's theorem (infinitely many primes): each
new prime $p$ contributes $J_2(p) = p^2 - 1 > 0$ to the sum.

\begin{correspondence}[Unconditional Infrared Divergence]
The divergence $\sigma(N) \to \infty$ is an \emph{infrared divergence}
of the Ramanujan scattering matrix: the total spectral weight
grows without bound. In the East Wing, this divergence is
\emph{desirable}---it ensures completeness. In the West Wing,
the same divergence is \emph{balanced} by the analytic corrections
in the Gram matrix, and whether this balance is perfect
(leading to $d_N^2 \to 0$) is exactly the content of RH.
\end{correspondence}

\subsection{The Mertens--Ramanujan Connection}

The Mertens--Ramanujan files (\lean{MertensRamanujan.lean},
\lean{MertensHarmony.lean}, \lean{BilinearMertens.lean})
establish the connection between the Mertens function (the
``magnetization'') and the Ramanujan matrix:
\begin{itemize}
  \item The \emph{geometric Mertens} identity
    (\lean{GeometricMertens.lean}) connects the Mertens product
    $\prod_{p \leq N}(1 - 1/p) \sim e^{-\gamma}/\ln N$ to the
    diagonal of the Ramanujan matrix.
  \item The \emph{bilinear Mertens} variance
    (\lean{BilinearMertens.lean}) bounds the off-diagonal
    Mertens interaction via the second moment $\sum \mu(j)\mu(k)/jk$.
  \item The \emph{log-correction} forms
    (\lean{LogCorrectionForm.lean}, \lean{LogCorrAsymptotics.lean})
    capture the subleading $O(1/\ln N)$ corrections to the
    Mertens-weighted quadratic form.
\end{itemize}

\begin{correspondence}[The Mertens Thermometer]
The Mertens product $e^{-\gamma}/\ln N$ is the \emph{temperature}
of the prime number gas at scale~$N$. The Ramanujan matrix entries
scale as $\gcd^2/(jk)$, which at temperature $T = 1/\ln N$ gives
an effective coupling constant $g_{\text{eff}} \sim \gcd^2/(jk\ln N)$.
As $N \to \infty$ ($T \to 0$), the coupling weakens---this is the
\emph{asymptotic freedom} of the Ramanujan interaction.
\end{correspondence}
